\documentclass[11pt]{article}
\usepackage[margin=1in]{geometry}
\usepackage[T1]{fontenc}
\usepackage{lmodern,amsmath,amssymb,amsthm,microtype}
\usepackage[hidelinks]{hyperref}
\newtheorem{theorem}{Theorem}
\newtheorem{lemma}{Lemma}
\newcommand{\E}{\mathbb E}
\newcommand{\R}{\mathbb R}
\newcommand{\supp}{\operatorname{supp}}
\newcommand{\Cov}{\operatorname{Cov}}
\newcommand{\RelInf}{\operatorname{RelInf}}
\title{Polynomial compatibility on the Boolean cube\\from disjoint-support separators}
\author{Alec Kriebel\\
\small\href{https://orcid.org/0009-0001-9320-500X}{ORCID: 0009-0001-9320-500X}}
\date{17 October 2026}
\begin{document}
\maketitle
\begin{abstract}
We give an explicit affirmative answer to the two-polynomial question posed by
Austrin in Oberwolfach Report 9/2025. Two nonzero real multilinear polynomials
of degree at most \(d\) on the uniform Boolean cube have a common nonzero
vertex whenever every relative coordinate influence is at most
\(1/(4096d^8)\). In fact, disjoint supports force a relative influence at least
\(1/(1024d^8)\). The result is a deduction from the disjoint-support separator
of Li--Li--Li--Liu and the covariance inequality of
O'Donnell--Saks--Schramm--Servedio, combined with a square-density estimate.
We supply the complete deduction and an explicit separator constant.
No novelty or first-priority claim is made for this consequence.
\end{abstract}

\section{The question and the quantitative answer}
All expectations refer to independent uniform signs \(X\in\{-1,1\}^n\).
For a nonzero real multilinear polynomial
\(a(x)=\sum_{S\subseteq[n]}\widehat a(S)\prod_{i\in S}x_i\), set
\[
 \|a\|_2^2=\E a(X)^2=\sum_S\widehat a(S)^2,\qquad
 \RelInf_i(a)=\frac{\sum_{S\ni i}\widehat a(S)^2}{\|a\|_2^2}.
\]
The denominator includes the constant coefficient. Austrin's question
\cite[p.~445]{OWR} asks whether an inverse-polynomial influence bound in the
degree guarantees a common nonzero vertex for two such polynomials.
The following gives that guarantee with absolute constants.

\begin{theorem}\label{thm:main}
Let \(n,d\ge1\) be integers and let \(f,g\) be nonzero real multilinear
polynomials on \(\{-1,1\}^n\), each of degree at most \(d\).
If \(f(x)g(x)=0\) at every vertex, then
\begin{equation}\label{eq:necessary}
 \max_{i\in[n]}\max\{\RelInf_i(f),\RelInf_i(g)\}
 \ \ge\ \frac{1}{1024d^8}.
\end{equation}
Consequently, if every relative influence of both polynomials is at most
\(1/(4096d^8)\), some vertex satisfies \(f(x)g(x)\ne0\).
\end{theorem}

The structural input is the prior separator of Li, Li, Li and Liu
\cite[Lemma~3.5]{LLLL}; the covariance input is already
\cite[\S3.3, Eq.~(6)]{OSSS}. Neither is a new theorem of this note.
The Markov/Bernoulli presentation in \cite[Fact~5 and Corollary~1]{KKDWY}
makes the numerical constant transparent. The resulting target-specific
application is written out below, including arbitrary common zeros.
The general distribution formulations in \cite[Conjectures~1.2 and~5.5]{ACC}
and other groups or input measures are outside the theorem asserted here.

\section{A quantitative total separator}
We reconstruct the prior separator to keep both its constant and its
behavior on common zeros explicit. Here alone we use \(\{0,1\}\)
coordinates; the affine change between the two cubes preserves degrees
and query depths.

\begin{lemma}[Li--Li--Li--Liu separator, with an extracted constant]
\label{lem:separator}
If real multilinear \(p,q\) of degree at most \(d\ge1\) on \(\{0,1\}^n\)
satisfy \(pq=0\) everywhere, there is a total deterministic decision tree
of depth at most \(16d^4\), with output \(h\in\{0,1\}\), such that
\(h=1\) on \(\supp p\) and \(h=0\) on \(\supp q\).
No coordinate repeats along a branch.
\end{lemma}
\begin{proof}
We first recall the block bound used in \cite{LLLL,KKDWY}.
Suppose \(R:\{0,1\}^b\to\R\) has degree at most \(m\), satisfies
\(|R|\le1\), \(R(0)=1\), and \(R(e_j)\le0\) for all \(j\).
Its expectation \(P(t)\) at independent Bernoulli(\(t\)) bits is a
polynomial of degree at most \(m\), bounded by one throughout \([0,1]\).
Moreover,
\[
 P'(0)=\sum_{j=1}^b(R(e_j)-1)\le-b.
\]
The classical Markov inequality, rescaled to \([0,1]\), gives
\(|P'(0)|\le2m^2\). Thus \(b\le2m^2\).

At a node, restrict \(p,q\) by the coordinates already queried.
If one restriction is zero, label for the other side, with an arbitrary
label if both are zero. Otherwise neither restriction can be a nonzero
constant, since the other has disjoint support.
Put \(D=\max(\deg p,\deg q)\ge1\). The multilinear representative of
\(s=p^2-q^2\) has degree at most \(2D\).
Choose a vertex \(z\) maximizing \(|s|\), whose maximum is positive.
Suppose \(p(z)\ne0\), so \(s(z)>0\) and \(q(z)=0\).

Choose an inclusion-maximal family of pairwise disjoint supports
\(M_1,\ldots,M_b\) of maximum-degree monomials of \(q\).
For each \(M_j\), fixing all variables outside it to their values at \(z\)
preserves its top coefficient: a further contribution would require a
strictly containing monomial. This restriction is therefore nonzero,
although it vanishes at \(z\). Some nonempty \(E_j\subseteq M_j\)
can be flipped from \(z\) to make \(q\ne0\).
There \(p=0\) and \(s<0\). The blocks \(E_j\) are disjoint.

Encode these block flips by \(b\) Boolean variables and divide \(s\)
by \(|s(z)|\). The resulting function \(R\) is bounded by one, has
\(R(0)=1\) and \(R(e_j)<0\), and has degree at most \(2D\):
each original coordinate becomes an affine function of one block bit,
and multilinearization cannot increase degree.
The block bound gives \(b\le8D^2\).

The union \(H=\bigcup_jM_j\) contains at most \(8d^3\) coordinates
and hits every maximum-degree monomial of \(q\), by maximality.
Query all coordinates in \(H\). On every answer branch, \(q\) becomes
zero or strictly loses degree; the degree of \(p\) cannot increase.
The case \(q(z)\ne0\) is symmetric.
Each nonterminal stage decreases \(\deg p+\deg q\) by at least one,
from an initial value at most \(2d\).
Thus at most \(2d\) stages, each costing at most \(8d^3\) queries,
suffice. Only free coordinates are queried.
Zero restrictions terminate before their degree is needed.
These rules give a total tree, including on the common-zero set.
\end{proof}

\section{Covariance and squared densities}
For a real function \(u\) on the sign cube, let
\[
 A_i(u)=\E|u(X)-u(X^{(i)})|,
\]
where \(X^{(i)}\) independently resamples coordinate \(i\);
it leaves that coordinate unchanged with probability one-half.

\begin{lemma}[OSSS covariance inequality]\label{lem:covariance}
Let a total deterministic tree, without repeated branch queries,
compute \(h:\{-1,1\}^n\to\{0,1\}\).
If \(r_i\) is the probability that it queries coordinate \(i\), then
for every real \(u\),
\begin{equation}\label{eq:covariance}
 |\Cov(h,u)|\le\frac12\sum_i r_i A_i(u).
\end{equation}
\end{lemma}
\begin{proof}
This is directly \cite[\S3.3, Eq.~(6)]{OSSS}, applied to \(2h-1\).
We give the finite-cube argument for completeness.
Run the tree solely on \(X\), and take an independent uniform \(Y\).
Let \(Z_t\) replace the first \(t\) queried coordinates of \(X\) by
those of \(Y\), and let \(L\) be the branch length.
A complete transcript is precisely a cylinder fixing the queried
\(X\)-coordinates. Conditional on it, \(Z_L\) is uniform:
queried coordinates come from \(Y\), and unqueried \(X\)-coordinates
remain independent fair signs. It is therefore independent of the
transcript and its output. Consequently
\[
 \Cov(h,u)=\E\left[(h(X)-\tfrac12)
             \sum_{t=1}^L\big(u(Z_{t-1})-u(Z_t)\big)\right].
\]
Conditional on a \emph{pre-query} transcript whose next index is \(i\),
the previously fixed \(X\)-coordinates have already been replaced
by independent \(Y\)-coordinates. All remaining \(X\)-coordinates
are unrestricted, and \(Y_i\) is fresh because queries do not repeat.
Thus \((Z_{t-1},Z_t)\) has exactly the unconditional uniform
coordinate-resampling-\(i\) law under this conditioning.
Using \(|h-\tfrac12|=\tfrac12\), taking the absolute telescoping sum,
and summing the probabilities of these pre-query cylinders gives
\eqref{eq:covariance}. All sums are finite, including variable-length
branches; \(u\) need not be computed by the tree.
\end{proof}

\begin{lemma}[Square-density estimate]\label{lem:square}
For any nonzero real \(a\), with
\(\rho_i=\RelInf_i(a)\) defined from its unique multilinear expansion,
\[
 A_i\!\left(\frac{a^2}{\|a\|_2^2}\right)
 \le2\sqrt{\rho_i(1-\rho_i)}\le2\sqrt{\rho_i}.
\]
\end{lemma}
\begin{proof}
Write \(a=b+x_ic\), with \(b,c\) independent of \(x_i\).
Parseval gives \(\|a\|_2^2=\E b^2+\E c^2\) and
\(\rho_i=\E c^2/\|a\|_2^2\).
A sign flip changes \(a^2\) by \(4x_ibc\), while independent resampling
flips with probability one-half. Hence the left side equals
\(2\E|bc|/\|a\|_2^2\).
Cauchy--Schwarz gives the asserted bound, including \(\rho_i=0,1\).
\end{proof}

\begin{proof}[Proof of Theorem~\ref{thm:main}]
Assume disjoint supports and take the total separator from
Lemma~\ref{lem:separator}, of depth \(T\le16d^4\).
Set \(u=f^2/\|f\|_2^2\), \(v=g^2/\|g\|_2^2\), and
\(\mu=\E h\). Then \(\E u=\E v=1\), \(hu=u\), and \(hv=0\), so
\[
 |\Cov(h,u)|+|\Cov(h,v)|=(1-\mu)+\mu=1.
\]
These identities also hold at common zeros, where both densities vanish.
With \(\delta\) the maximum relative influence of either polynomial,
Lemmas~\ref{lem:covariance}--\ref{lem:square} give
\[
 1\le \sum_i r_i\!
 \left(\sqrt{\RelInf_i(f)}+\sqrt{\RelInf_i(g)}\right)
 \le2\Big(\sum_i r_i\Big)\sqrt\delta
 \le2T\sqrt\delta\le32d^4\sqrt\delta.
\]
This proves \eqref{eq:necessary}.
If \(\delta\le1/(4096d^8)\), the last upper bound is at most \(1/2\),
a contradiction.
\end{proof}

A nonzero constant has full support, so constants cause no exception.
The proof also covers \(n<d\).
The earlier exponential sufficient bound uses
\(\delta<2^{-d}/d\), as in \cite[\S5.2]{ACC}.
Its non-strict endpoint printed in \cite{OWR} fails already for
\(d=1\), \(f=1+x_1\), \(g=1-x_1\), both of relative influence \(1/2\).
The smaller inclusive threshold in Theorem~\ref{thm:main} has slack.
No optimality of its exponent or constants is asserted.

\section{Attribution, verification, and review status}
The contribution of this note is the explicit application above to the
stated two-polynomial question. The separator, the real covariance
inequality, and the Markov/Bernoulli method are credited prior inputs,
available before this note. Reconstructing their proofs does not establish
novelty. A bounded primary-source audit found no earlier recorded complete
target-specific deduction in the bodies inspected; it does not certify
historical priority. In particular, the full ePrint version of the 2022
paper and the full 2023 non-interactive-commitment followup were
inaccessible to that audit. Detailed source/version records, inspection
limits, and independent priority reviews accompany the note.
No firstness or current global literature-status claim is made.

The accompanying source and verification archive includes the original
proof record, new independent mathematical audits, exact rational
diagnostics, and a portable offline runner. Finite enumerations check
adaptive trees, support geometry, Fourier identities, boundary cases,
and deliberate incorrect variants in ordinary and optimized execution.
They supplement the universal written proof; they are not a formal
proof certificate or an exhaustive check in arbitrary dimension.
Third-party source bodies are not redistributed.

AI tools were used extensively in developing, writing, reproducing and
adversarially reviewing this work. The reviews are internal AI-assisted
checks. This preprint is unrefereed; no external human peer review,
journal acceptance, or proof-assistant certification is claimed.

\begin{thebibliography}{9}
\bibitem{OWR}
P.~Austrin, \emph{Zeros of Low-Degree, Low-Influence Polynomials on the
Boolean Hypercube}, contributed problem in
MFO--RIMS Tandem Workshop: Optimization, Theoretical Computer Science
and Algebraic Geometry: Convexity and Beyond,
Oberwolfach Reports \textbf{22} (2025), no.~1, 405--450, at p.~445.
\href{https://doi.org/10.4171/OWR/2025/9}{doi:10.4171/OWR/2025/9}.

\bibitem{ACC}
P.~Austrin, H.~Chung, K.-M.~Chung, S.~Fu, Y.-T.~Lin and M.~Mahmoody,
\emph{On the Impossibility of Key Agreements from Quantum Random Oracles},
CRYPTO 2022, LNCS \textbf{13508}, 165--194.
\href{https://doi.org/10.1007/978-3-031-15979-4_6}{doi:10.1007/978-3-031-15979-4\_6}.

\bibitem{LLLL}
L.~Li, Q.~Li, X.~Li and Q.~Liu,
\emph{Impossibility of Perfectly Complete Many-Round Key Agreement
in the QROM}, 4 August 2026.
\href{https://arxiv.org/abs/2608.03824v1}{arXiv:2608.03824v1},
Lemma~3.5, pp.~7--8.

\bibitem{KKDWY}
R.~Kothari, M.~Kovacs-Deak, D.~Wang and R.~Z.~Yang,
\emph{Rational degree is polynomially related to degree},
28 September 2026.
\href{https://arxiv.org/abs/2601.08727v3}{arXiv:2601.08727v3},
Fact~5, Theorem~1 and Corollary~1.

\bibitem{OSSS}
R.~O'Donnell, M.~Saks, O.~Schramm and R.~A.~Servedio,
\emph{Every decision tree has an influential variable},
FOCS 2005, 31--39.
\href{https://doi.org/10.1109/SFCS.2005.34}{doi:10.1109/SFCS.2005.34};
\href{https://arxiv.org/abs/cs/0508071v1}{arXiv:cs/0508071v1},
\S3.3, Eq.~(6), p.~8.
\end{thebibliography}
\end{document}
